\section{A specified rational field with 783 hyperbolic cycles}\label{app:783-proof}
At \(N=32\), Proposition~\paperlink{main}{theorem.6.1}{prop:stationary-budget}{6.1} of the article gives 783 cycles at degree 30 for sufficiently separated parameters. Here we fix the parameters in the unreduced family of Theorem~\ref{thm:stationary-turn} and verify its returns at degree 31.

\begin{proposition}[Explicitly specified $783$-cycle rational vector field]\label{prop:783-witness}
Fix $(n,d)=(15,15)$, $r=29$, $B=96$, $q=2^{-64}$, $A=2^{10{,}000{,}000}$, and $\lambda=2^{-20{,}000{,}000}$. Let $G(y)$ be the explicit rational even polynomial~\eqref{eq:st-choice} of degree $30$, $F(y)=\int_0^y G(t)\dd t$, $U(x)=A^2 T_{15}(x)^2$, and
\[
  P(x,y) = 31y^{30} - 29y^{28}, \qquad
  Q(x,y) = -U'(x) + A\lambda F(y).
\]
Then $X=(P,Q)\in\Q[x,y]^2$ is coprime, has exact component supports $\splitcount{2}{30}$ ($N(X)=32$) and actual degree $D(X)=31$, and possesses at least $783=29\times(14+13)$ distinct simultaneous hyperbolic limit cycles. In each of the $29$ disjoint horizontal strips, the local return section is parametrized by $h=X^2+V(Y)$. At least one cycle meets each of the $14$ return intervals in $h>0$ and each of the $13$ in $h<0$.
\end{proposition}

In the coordinates of \eqref{eq:st-chart}, the physical Hamiltonian \(H=U(x)+V(y)\) equals \(h\) on a minimum strip and \(A^2-h\) on a maximum strip. For the perturbed return, \(h\) specifies the starting point on the section.

We prove Proposition~\ref{prop:783-witness} at the fixed rational parameters $n=d=15$, $r=29$, $B=96$, $q=2^{-64}$, $A=2^{10^7}$, and $\lambda=2^{-2\cdot 10^7}$, distinguishing explicitly between finite rational/interval enclosures verified computationally and uniform complex-analytic estimates holding across all ovals in the $27$ brackets.

\paragraph{1. Finite rational and outward-rounded interval enclosures.}
Here $\tilde{e}_0,\dots,\tilde{e}_{14}$ equal $0,31,35,39,43,47,51,55,58,59,63,67,71,75,79$. Exact rational polynomial arithmetic reconstructs $D_0,\dots,D_{14}\in\Q[y]$, $G(y)$, $F(y)$, and $-U'(x)$, confirming that $P$ has $2$ nonzero terms, $Q$ has $15$ odd $x$-powers and $15$ odd $y$-powers ($N(X)=32$, $D(X)=31$, $\gcd(P,Q)=1$), $\sum_k|[y^k]D_i|<2^{40}$ ($0\le i\le 14$), and $\sum_k|[y^k]F|<2^{60}$. Outward-rounded Arb ball arithmetic verifies $2^{-12}<|c_{i,\pm}|<2^{12}$ ($0\le i\le 14$) and $c_{i,+}c_{i,-}<0$ ($1\le i\le 14$). For each of the $14$ sign-exchanging adjacent pairs on $h>0$ and $13$ pairs on $h<0$ ($i\ge 1$), set $\Delta:=\tilde{e}_{i+1}-\tilde{e}_i\ge 1$, $t_0:=|c_{i,\pm}/c_{i+1,\pm}|^{1/\Delta}$, $w_0:=q^{2B-\tilde{e}_i-\tilde{e}_{i+1}}t_0$, $w=w_0 v$, and $N_{\mathrm{scale}}(v):=\sigma_i c_{i,\pm}q^{(B-\tilde{e}_i)^2}(w_0 v)^{\tilde{e}_i}$. Define explicit dyadic rational endpoints $w_-:=2^{-80}q^{2B-\tilde{e}_i-\tilde{e}_{i+1}}\lfloor 2^{80}t_0(1-\frac{1}{5\Delta})\rfloor$ and $w_+:=2^{-80}q^{2B-\tilde{e}_i-\tilde{e}_{i+1}}\lceil 2^{80}t_0(1+\frac{1}{5\Delta})\rceil$, and put $[v_-,v_+]:=[w_-/w_0,w_+/w_0]$. Outward-rounded Arb evaluation certifies that on all $27$ brackets: (i) $[v_-,v_+]\subset(1-\frac{1}{4\Delta},1+\frac{1}{4\Delta})$ and the $w$-intervals on each side are pairwise disjoint inside $(2^{-10998},2^{-2002})$, so $H_*:=2^{-10^6}<|h|<2^{-100}$ and $|N_{\mathrm{scale}}(v)|>2^{-2\cdot 10^6}$; and (ii) the normalized leading-power sum $L_{\mathrm{lead}}(v):=\sum_{k=0}^{14}\frac{\sigma_k c_{k,\pm}}{\sigma_i c_{i,\pm}}q^{(B-\tilde{e}_k)^2-(B-\tilde{e}_i)^2}(w_0 v)^{\tilde{e}_k-\tilde{e}_i}$ satisfies $L_{\mathrm{lead}}(v_-)>+\frac{1}{8}$, $L_{\mathrm{lead}}(v_+)<-\frac{1}{8}$, and $L_{\mathrm{lead}}'(v)<-\frac{\Delta}{2}$ throughout $[v_-,v_+]$.

\paragraph{2. Uniform holomorphic moment remainder on $|w|\le 1/8$ (analytic proof).}
Write $\Psi_i^\pm(w):=I[D_i](\pm w^{58})$. On $|y|=3/4$, $|y(1-y^2)^{1/29}|\ge\frac{3}{4}(\frac{7}{16})^{1/29}>\frac{1}{2}$, so Rouch\'e's theorem yields an analytic left-turn inverse $y=\psi(z)$ on $|z|\le 1/2$ with $|\psi(z)|\le 3/4$ and $\psi(1/4)\in(1/5,1/3)$. Cauchy's estimate on $|z|\le 3/8$ gives $|\psi'(z)|\le\frac{3/4}{1/2-3/8}=6\le 8$ and $|[z^{2m}]\psi^{2j}\psi'|\le 8(8/3)^{2m}$. Taking $z_0=1/4$, for $|w|\le 1/4$ ($|h|\le 2^{-116}$, $|h|/z_0^{29}\le 2^{-58}$), the series~\eqref{eq:st-upper-series} is bounded by $4z_0^{2m+31/2}$, while $|b_m|\le\frac{1}{2m+1}$ and $|\cos(\pi\frac{2m+1}{29})|\ge\sin\frac{\pi}{58}>\frac{1}{29}$ bound the fractional constants by $30$, so the left-cut integral is bounded by $\sum_{m\ge 0}8(8/3)^{2m}(4\cdot 4^{-2m-31/2}+30\cdot 4^{-31-4m})<1$. On the middle real segment $[\psi(1/4),v_0]$ with $v_0:=1-2^{-16}$, the unimodal function $-V(y)=y^{29}(1-y^2)$ exceeds $\min(-V(1/5),-V(v_0))>2^{-75}\gg 2^{-116}\ge|h|$, giving a holomorphic square root of modulus $<2$. Near the right turn $y=1$, $V(1)=0$, $V'(1)=2$, and $|V''(y)|<2^{12}$ on $|y-1|\le 2^{-12}$, so on $|y-1|=2^{-12}$ and $|v|\le 2^{-13}$ one has $|V(y)-2(y-1)|+|v|\le 2^{-13}+2^{-13}<2^{-11}=|2(y-1)|$; Rouch\'e's theorem therefore yields a holomorphic right-turn inverse $\phi_{\mathrm{right}}(v)$ on $|v|\le 2^{-13}$ with $|\phi_{\mathrm{right}}'(v)|\le 1$. Since $2^{-17}<-V(v_0)<2^{-14}$, parametrizing $[v_0,u(h)]$ by $v=V(v_0)+t(h-V(v_0))$ ($0\le t\le 1$) bounds $|I_j(\pm w^{58})|\le 2^{12}$ on $|w|\le 1/4$. Multiplying by $\sum_k|[y^k]D_i|<2^{40}$ gives $|\Psi_i^\pm(w)|\le 2^{52}<2^{100}$ on $|w|\le 1/4$, and Cauchy's estimate on $|w|\le 1/8$ ($\tilde{e}_i\le 79$) yields
\begin{equation}\label{eq:app-psi-rem}
  \bigl|\Psi_i^\pm(w)/w^{\tilde{e}_i} - c_{i,\pm}\bigr| \le 2^{100}\cdot 4^{79}\cdot 8|w| < 2^{270}|w| \qquad (|w|\le 1/8).
\end{equation}
On the complex disk $|v-1|\le\frac{1}{2\Delta}$ ($2^{-11000}<|w|<2^{-2000}$),~\eqref{eq:app-psi-rem} bounds the difference $|I[G](\pm w^{58})/N_{\mathrm{scale}}(v)-L_{\mathrm{lead}}(v)|$ by $2^{-1700}$ (and its $v$-derivative on $|v-1|\le\frac{1}{4\Delta}$ by $2^{-1690}$ via Cauchy), so $I[G]/N_{\mathrm{scale}}$ inherits strict margins $>+\frac{1}{8}-2^{-1700}$, $<-\frac{1}{8}+2^{-1700}$, and slope $<-\frac{\Delta}{2}+2^{-1690}$ on every $[v_-,v_+]$.

\paragraph{3. Complex angular ODE and full nonlinear return (analytic proof).}
In clockwise polar coordinates $(X,Y)=(\xi\cos\phi,a-\xi\sin\phi)$ around $(0,a)$, on every reference oval ($H_*<|h|<2^{-100}$) one has $2^{-16}<\xi<4$, $|X|<2$, $-1/4<Y<5/4$, and angular speed numerator $J:=2X^2+(Y-a)V'(Y)=2X^2+31Y^{28}(Y-a)^2(Y+a)>2^{-40}H_*$. For $f(\xi):=X^2+V(Y)$, $f'(\xi_0)=J/\xi_0>2^{-43}H_*$ and $|f''|<2^{100}$ on $|\delta\xi|\le 1/4$, so Rouch\'e's theorem on $|\delta\xi|=\beta:=2^{-160}H_*$ yields a holomorphic implicit radius $\xi(\phi,h')$ for $|h'-h|\le\eta:=2^{-220}H_*^2$ with $|J|>2^{-62}H_*$, $|X|,|Y|<3$, $|g_i|<2$, $|g_i-c_i|<2^6/A$, $|V'|<2^{56}$, and $|F|<2^{111}$. Along the perturbed angular ODE
\[
  \frac{\dd h}{\dd\phi} = \frac{\lambda\xi^2 V'(Y)g_i(X)F(Y)}{J - \lambda X g_i(X)F(Y)},
\]
for $|\lambda|\le R:=2^{-600}H_*^3$ we have $|\dd h/\dd\phi|\le 2^{250}|\lambda|/H_*$ and $2\pi\cdot 2^{250}R/H_*<\eta/4$, so starting at $|h'-h|\le\eta/2$ yields a holomorphic one-turn displacement $\mathcal{D}(h',\lambda)$ of modulus $<1$. Cauchy's tail bound in $\lambda$ ($2\lambda/R^2$), replacement of $g_i$ by $c_i$ ($2^{260}/(A H_*)$), and Cauchy's $h'$-derivative factor ($4/\eta$) give, at $\lambda=2^{-2\cdot 10^7}$ and $A=2^{10^7}$,
\[
  \left|\frac{\mathcal{D}(h,\lambda)}{2c_i\lambda} - I[G](h)\right| < 2^{-3{,}000{,}000}, \qquad
  \left|\partial_h\!\left(\frac{\mathcal{D}(h,\lambda)}{2c_i\lambda} - I[G](h)\right)\right| < 2^{-3{,}000{,}000}.
\]
Dividing by $|N_{\mathrm{scale}}(v)|>2^{-2\cdot 10^6}$ leaves normalized value error $<2^{-10^6}$. Since $|\dd h/\dd v|=58|h/v|<120$ and $|N'_{\mathrm{scale}}/N_{\mathrm{scale}}|=\tilde e_i/|v|<106$, the product rule bounds the normalized slope error by $(120+106)2^{-10^6}<2^{-100}$ on all $27$ brackets in each of the $29$ disjoint horizontal strips, certifying $29\times 27=783$ simultaneous hyperbolic limit cycles.

\paragraph{4. Supplementary whole-oval quadrature check and rigorous right-turn root enclosure.}
As an independent numerical cross-check of Lemma~\ref{lem:st-expansion} on both sides of $h=0$, we evaluate $I[D_i](h)$ directly by $1000$-bit outward-rounded Arb integration at $\ell=\pm 2^{-16}$ ($h=V(\ell)$). Writing $u(h)=1+z_*$ for the unique root of $V(y)=h$ in $(a,\infty)$, exact rational evaluation first verifies endpoint sign-straddling $V(1-|h|)<h<V(1+|h|)$ together with strict monotonicity $V'(1+I_0)>1$ across $I_0:=[-|h|,|h|]$. For $T(z):=h/((1+z)^{29}(2+z))$, ball arithmetic verifies strict self-map inclusion $T(I_0)\subset\operatorname{int}(I_0)$, so every iterate $I_m:=T^m(I_0)$ encloses $z_*$ by inclusion monotonicity ($z_*=T(z_*)\in I_m$); moreover, on the final narrow enclosure $[u_{\mathrm{lo}},u_{\mathrm{hi}}]$ around $1+I_5$, interval evaluation also verifies $V'([u_{\mathrm{lo}},u_{\mathrm{hi}}])>1$ on the hull and strict opposite residual signs $V(u_{\mathrm{lo}})-h<0<V(u_{\mathrm{hi}})-h$ at its endpoints before factoring $h-V(y)=(y-\ell)(u(h)-y)B(y)$.

